%% LyX 2.3.8 created this file.  For more info, see http://www.lyx.org/.
%% Do not edit unless you really know what you are doing.
\documentclass[english]{IEEEConf}
\usepackage[T1]{fontenc}
\usepackage[latin9]{inputenc}
\usepackage{geometry}
\geometry{verbose,tmargin=0.5in,bmargin=0.5in,lmargin=0.5in,rmargin=0.5in}
\usepackage{color}
\usepackage{amsmath}
\usepackage{amsthm}
\usepackage{amssymb}
\usepackage{esint}

\makeatletter
%%%%%%%%%%%%%%%%%%%%%%%%%%%%%% Textclass specific LaTeX commands.
\theoremstyle{definition}
\newtheorem{defn}{\protect\definitionname}
\theoremstyle{plain}
\newtheorem{thm}{\protect\theoremname}
\theoremstyle{definition}
\newtheorem{assumption}{Assumption}
\theoremstyle{remark}
\newtheorem{rem}{\protect\remarkname}
\theoremstyle{definition}
\newtheorem{property}{Property}

\makeatother

\usepackage{babel}
\providecommand{\definitionname}{Definition}
\providecommand{\remarkname}{Remark}
\providecommand{\theoremname}{Theorem}

\begin{document}
\title{\vspace{0.25in}
}
\title{Exponentially Convergent Modular Adaptive Control for Nonlinear Target
Tracking}
\author{Max L. Gardenswartz, Cristian F. Nino, Scott A. Nivison, Warren E.
Dixon\thanks{Max L. Gardenswartz, Brandon C. Fallin, and Warren E. Dixon are with
the Department of Mechanical and Aerospace Engineering, University
of Florida, Gainesville, FL 32611, USA. Email: \{mgardenswartz, brandonfallin,
wdixon\}@ufl.edu. Cristian F. Nino is with the Florida Institute for
Human and Machine Cognition, Pensacola, FL 32502, USA. Email: cnino@ihmc.org.}\thanks{This research is supported in part by AFRL project FA8651-24-1-0018
and AFOSR grant FA9550-19-1-0169. Any opinions, findings and conclusions
or recommendations expressed in this material are those of the author(s)
and do not necessarily reflect the views of the sponsoring agency.}}
\maketitle
\begin{abstract}
A
\end{abstract}

\begin{IEEEkeywords}
Neural networks, Adaptive control, Stability of nonlinear systems
\end{IEEEkeywords}


\section{Introduction}

Depth with nonlinearity creates no bad local minima in ResNets \cite{Kawaguchi.Bengio2019}.
\textcolor{red}{CHANGE ALL C2 to LIPSCHITZ. YOUR EPSILON ARGUMENT
ARE WRONG.}

\section{Preliminaries}

\subsection{Notation}

Let $n\in\mathbb{Z}_{\ge1}$. The $\left\Vert \right\Vert $ means
2-norm. A continuous function $\mathcal{C}^{2}$...Zero vector. $\subsetneq$.
Ball or radius $r$. Closure. Filippov regularization. Signum. $C^{1}$
norm.

\subsection{Nonsmooth Analysis }

See latest Cristian paper.

\subsection{\label{subsec:Deep-Residual-Neural}Deep Residual Neural Network
Model}

\textcolor{blue}{The residual neural network (ResNet) architecture,
characterized by skip connections and hierarchical feature extraction,
models incremental changes rather than complete transformations of
the underlying nonlinear mapping. This architecture learns the differences
(or ``residuals'') between the input and desired output at each
layer, thereby enabling effective function approximation for complex
nonlinear systems without requiring explicit governing equations.
A fully-connected feedforward ResNet is constructed as follows.}

\textcolor{blue}{Let} $b\in\mathbb{Z}_{\ge0}$ \textcolor{blue}{be
the number of }building blocks, \textcolor{blue}{let the input be}
$x\in\mathbb{R}^{L_{\text{in}}}$, and \textcolor{blue}{let the output
be} $y\in\mathbb{R}^{L_{\text{out}}}$, \textcolor{blue}{where $L_{\text{in}},L_{\text{out}}\in\mathbb{Z}_{\ge1}$}.
For each block index $i\in\left\{ 0,\ldots,b\right\} $, let $k_{i}\in\mathbb{Z}_{\ge1}$
denote the number of hidden layers in the $i^{\text{th}}$ block,
let $\kappa_{i}\in\mathbb{R}^{L_{i,0}}$ denote the block input, and
let $\theta_{i}\in\mathbb{R}^{p_{i}}$ denote the vector of parameters
(weights and biases) associated with the $i^{\text{th}}$ block. \textcolor{blue}{Note
that $\kappa_{0}\triangleq x$ and $L_{0,0}\triangleq L_{\text{in}}$.}

For all $\left(i,j\right)\in\left\{ 0,\ldots,b\right\} \times\left\{ 0,\ldots,k_{i}+1\right\} $,
let $L_{i,j}\in\mathbb{Z}_{\ge1}$ denote the number of neurons in
the $j^{\text{th}}$ layer for the $i^{\text{th}}$ block. For all
$\left(i,j\right)\in\left\{ 0,\ldots,b\right\} \times\left\{ 0,\ldots,k_{i}\right\} $,
define the augmented dimension $L_{i,j}^{a}\triangleq L_{i,j}+1$.
Each block function $\Phi_{i}:\mathbb{R}^{L_{i,0}^{a}}\times\mathbb{R}^{p_{i}}\to\mathbb{R}^{L_{i,k_{i}+1}}$
is a fully connected, feedforward \textcolor{blue}{deep neural network}
(DNN) with \textcolor{blue}{$L_{i,k_{i}+1}\triangleq L_{\text{out}}$
for all $i\in\left\{ 0,\ldots,b\right\} $.} For any input $v\in\mathbb{R}^{L_{i,j}^{a}}$,
the DNN is defined recursively by

\begin{equation}
\varphi_{i,j}\left(v\right)\triangleq\begin{cases}
V_{i,0}^{\top}v, & j=0,\\
V_{i,j}^{\top}\phi_{i,j}\left(\varphi_{i,j-1}\left(v\right)\right), & j\in\left\{ 1,\ldots,k_{i}\right\} ,
\end{cases}\label{eq:DNN}
\end{equation}
with $\Phi_{i}\left(v,\theta_{i}\right)=\varphi_{i,k_{i}}\left(v\right)$.

For each $j\in\left\{ 0,1,\ldots,k_{i}\right\} $, the matrix $V_{i,j}\in\mathbb{R}^{L_{i,j}^{a}\times L_{i,j+1}}$
contains the weights and biases; in particular, if a layer has $n$
inputs \textcolor{blue}{(including the appended bias) }and the subsequent
layer has $m$ nodes, then $V\in\mathbb{R}^{n\times m}$ is constructed
so that its \textcolor{blue}{$\left(r,c\right)^{\text{th}}$ }entry
represents the weight from the \textcolor{blue}{$r^{\text{th}}$ }node
of the input to the \textcolor{blue}{$c^{\text{th}}$} node of the
output, with the last row corresponding to the bias terms. For the
DNN architecture described by (\ref{eq:DNN}), the vector of DNN weights
of the $i^{\text{th}}$ block is $\theta_{i}\triangleq\left[\begin{array}{ccc}
\text{vec}\left(V_{i,0}\right)^{\top} & \cdots & \text{vec}\left(V_{i,k_{i}}\right)^{\top}\end{array}\right]^{\top}\in\mathbb{R}^{p_{i}}$, where $p_{i}\triangleq\sum_{j=0}^{k_{i}}L_{i,j}^{a}L_{i,j+1}$,
and $\text{vec}\left(V_{i,j}\right)$ denotes the vectorization of
$V_{i,j}$ performed in column-major order (i.e., the columns are
stacked sequentially to form a vector). \textcolor{blue}{For all $\left(i,j\right)\in\left\{ 0,\ldots,b\right\} \times\left\{ 0,\ldots,k_{i}\right\} $,}
the \textcolor{blue}{ensemble} activation function \textcolor{blue}{of
the $j^{\text{th}}$ layer of the $i^{\text{th}}$ block is given
by $\phi_{i,j}:\mathbb{R}^{L_{i,j}}\to\mathbb{R}^{L_{i,j}^{a}}$,
where} \textcolor{blue}{$\phi_{i,j}\left(\varphi_{i,j-1}\right)$
$\triangleq$ $\left[\varsigma_{i,j}\left(\left(\varphi_{i,j-1}\right)_{1}\right)\right.$
$\begin{array}{c}
\varsigma_{i,j}\left(\left(\varphi_{i,j-1}\right)_{2}\right)\end{array}$ $\begin{array}{c}
\cdots\end{array}$ $\begin{array}{c}
\varsigma_{i,j}\left(\left(\varphi_{i,j-1}\right)_{L_{i,j}}\right)\end{array}$ $\left.1\right]^{\top}$} $\in\mathbb{R}^{L_{i,j}^{a}}$\textcolor{blue}{,
the activation function is} $\varsigma_{i,j}:\mathbb{R}\to\mathbb{R}$,
and 1 accounts for the bias term.

A pre-activation design is used so that, before each block (except
block $0$), the output of the previous block is processed by an external
activation function. Specifically, for each block $i\in\left\{ 1,\ldots,b\right\} $,
define the pre-activation mapping $\psi_{i}:\mathbb{R}^{L_{i,k_{i}+1}}\to\mathbb{R}^{L_{i,k_{i}+1}^{a}}$
by $\psi_{i}\left(\kappa_{i}\right)$ $=$ \textcolor{blue}{$\left[\varrho_{i}\left(\left(\kappa_{i}\right)_{1}\right)\right.$
$\varrho_{i}\left(\left(\kappa_{i}\right)_{2}\right)$ $\begin{array}{c}
\cdots\end{array}$ $\varrho_{i}\left(\left(\kappa_{i}\right)_{L_{i,k_{i}+1}}\right)$
$\left.1\right]^{\top}$ $\in\mathbb{R}^{L_{i,k_{i}+1}^{a}}$} where
$\left(\kappa_{i}\right)_{\ell}$ denotes the $\ell^{\text{th}}$
component of $\kappa_{i}$, the activation function is $\varrho_{i}:\mathbb{R}\to\mathbb{R}$,
and 1 accounts for the bias term. The output of $\psi_{i}$ serves
as the input to block $i$ and the residual connection is implemented
by adding the current block output to the pre-activated output from
the previous block. Hence, the ResNet recursion is defined by
\begin{equation}
\kappa_{i+1}\triangleq\begin{cases}
\Phi_{0}\left(\kappa_{0}^{a},\theta_{0}\right), & i=0,\\
\kappa_{i}+\Phi_{i}\left(\psi_{i}\left(\kappa_{i}\right),\theta_{i}\right), & i\in\left\{ 1,\ldots,b\right\} ,
\end{cases}\label{eq:ResNet}
\end{equation}
with output $y\in\mathbb{R}^{L_{\text{out}}}$ and overall parameter
vector $\Theta\triangleq\left[\begin{array}{ccc}
\theta_{0}^{\top} & \cdots & \theta_{b}^{\top}\end{array}\right]^{\top}\in\mathbb{R}^{{\tt p}}$, with ${\tt p}\triangleq\sum_{i=0}^{b}p_{i}$, where $\kappa_{0}^{a}\triangleq\left[\begin{array}{cc}
\kappa_{0}^{\top} & 1\end{array}\right]^{\top}\in\mathbb{R}^{L_{0,0}^{a}}$ denotes the augmented input to block $i=0$. Therefore, the complete
ResNet is represented as $\Psi:\mathbb{R}^{L_{\text{in}}}\times\mathbb{R}^{{\tt p}}\to\mathbb{R}^{L_{\text{out}}}$
expressed as $\Psi\left({\color{blue}x},\Theta\right)=\kappa_{b+1}$.

In order for ResNets to act as are universal function approximators
of continuous functions on compact domains, their activation functions
must be selected to belong to a class of functions called $\mathcal{A}_{\text{quad}}$
.
\begin{defn}
\emph{\label{def:a-quad}\cite[Definition 3.1]{Tabuada.Gharesifard2023}
The activation function $\varsigma:\mathbb{R}\to\mathbb{R}$ is said
to be of class $\mathcal{A}_{\text{quad}}$ (denoted $\varsigma\in\mathcal{A}_{\text{quad}}$)
if it is (i) nondecreasing, (ii) non-constant, and (iii) there exists
an $\ell\in\mathbb{Z}_{\ge0}$ such that $\varsigma$ is a globally-defined
solution to $\frac{d\zeta}{ds}=a_{0}+a_{1}\zeta+a_{2}\zeta^{2}$,
where $\zeta\triangleq\frac{d^{(\ell)}\varsigma}{ds^{(\ell)}}$ is
injective, $a_{0},a_{1},a_{2}\in\mathbb{R}$, $a_{2}\ne0$, and $\frac{d^{(0)}\varsigma}{ds^{(0)}}=\varsigma$.}\footnote{By consequence of Definition \ref{def:a-quad}, any $\varsigma\in\mathcal{A}_{\text{quad}}$
is necessarily an analytical function. Thus, $\varsigma$ and all
its derivatives are locally Lipschitz on $\mathbb{R}$. Hyperbolic
tangent, sigmoid, and softplus are of class $\mathcal{A}_{\text{quad}}$
but ReLU, swish, and the identity mapping are not. The function $\varsigma\left(s\right)=\frac{1}{c}\log\left(1+{\rm e}^{cs}\right)$
approximates ReLU for large $c\in\mathbb{R}_{>0}$ and is of class
$\mathcal{A}_{\text{quad}}$. Leaky ReLU is given as $\varsigma\left(s\right)=s$
for all $s\in\mathbb{R}_{\ge0}$ and $\varsigma\left(s\right)=rs$
(with $r\in\mathbb{R}_{>0}$) for all $s<0$. The function $\varsigma\left(s\right)=rs+\frac{1}{c}\log\left(1+{\rm e}^{\left(1-r\right)cs}\right)$
approximates leaky ReLU for large $c\in\mathbb{R}_{>0}$ and is of
class $\mathcal{A}_{\text{quad}}$.}
\end{defn}
%
\begin{thm}
\label{thm:Universal-Approximation-Theorem}\cite[Corollary 5.2]{Tabuada.Gharesifard2023}
(Universal Approximation Theorem for Residual Neural Networks). For
any compact domain $\Omega\subsetneq\mathbb{R}^{L_{\text{in}}}$,
continuous function $F:\Omega\to\mathbb{R}^{L_{\text{out}}}$, and
approximation accuracy $\overline{\varepsilon}\in\mathbb{R}_{>0}$,
there exists a number of blocks $b\in\mathbb{Z}_{>0}$ (thus, a ${\tt p}\in\mathbb{Z}_{>0}$)
and a parameter vector $\theta^{\ast}\in\mathbb{R}^{{\tt p}}$ such
that for a ResNet architecture $\Psi:\mathbb{R}^{L_{\text{in}}}\times\mathbb{R}^{{\tt p}}\to\mathbb{R}^{L_{\text{out}}}$
with blocks of width $w=\max\left(L_{\text{in}},L_{\text{out}}\right)+1$
(i.e., $L_{i,j}\ge w$ for all $\left(i,j\right)\in\left\{ 0,\ldots,b\right\} \times\left\{ 1,\ldots,k_{i}\right\} $)
and depth 1 (i.e., $k_{i}\ge1$ for all $i\in\left\{ 0,\ldots,b\right\} $),
\begin{equation}
\sup_{x\in\Omega}\left\Vert \varepsilon\left(x\right)\right\Vert \le\overline{\varepsilon},\label{eq:lemma-ufat-1}
\end{equation}
where $\varepsilon\left(x\right)\triangleq h\left(x\right)-\Psi\left(x,\theta^{\ast}\right)$,
provided that for all $\left(i,j\right)\in\left\{ 0,\ldots,b\right\} \times\left\{ 0,\ldots,k_{i}\right\} $,
the activation functions $\varsigma_{i,j}\in\mathcal{A}_{\text{quad}}$\emph{.}
\end{thm}
\textcolor{red}{HAVE TO MENTION HIGHER-ORDER APPROXIMATION.}

\section{Problem Formulation}

Consider the second-order nonlinear dynamical system described by
the differential equation,
\begin{equation}
\ddot{q}\left(t\right)=f\left(q\left(t\right),\dot{q}\left(t\right)\right)+u\left(t\right)+d\left(t\right),\label{eq:agent-dynamics}
\end{equation}
where $t\in\mathbb{R}_{\ge0}$ denotes time, $u:\mathbb{R}_{\ge0}\to\mathbb{R}^{n}$
denotes the control input, $f:\mathbb{R}^{n}\times\mathbb{R}^{n}\to\mathbb{R}^{n}$
denotes unknown drift dynamics of class $\mathcal{C}^{2}$, and $d:\mathbb{R}_{\ge0}\to\mathbb{R}^{n}$
denotes an unknown, time-variant exogenous disturbance of class $\mathcal{C}^{2}$.
For an initial condition set $\left(t_{0},q_{0},\dot{q}_{0}\right)\in\mathbb{R}_{\ge0}\times\mathbb{R}^{n}\times\mathbb{R}^{n}$,
let $q:\mathcal{I}\to\mathbb{R}^{n}$, $\dot{q}:\mathcal{I}\to\mathbb{R}^{n}$,
and $\ddot{q}:\mathcal{I}\to\mathbb{R}^{n}$ denote the generalized
position, velocity, and acceleration, respectively, of the resulting
solution to (\ref{eq:agent-dynamics}) such that $q\left(t_{0}\right)=q_{0}$
and $\dot{q}\left(t_{0}\right)=\dot{q}_{0}$ and whose interval of
existence is $\mathcal{I}\triangleq\left[0,T_{\max}\right)$ for $T_{\max}\in\mathbb{R}_{>t_{0}}\cup\left\{ \infty\right\} $.

The desired trajectory for (\ref{eq:agent-dynamics}) is governed
by the autonomous second-order system,
\begin{equation}
\ddot{q}_{d}\left(t\right)=h\left(q_{d}\left(t\right),\dot{q}_{d}\left(t\right)\right),\label{eq:reference-dynamics}
\end{equation}
where $h:\mathbb{R}^{n}\times\mathbb{R}^{n}\to\mathbb{R}^{n}$ is
an unknown function of class $\mathcal{C}^{2}$. For an initial condition
set $\left(t_{0},q_{d,0},\dot{q}_{d,0}\right)\in\mathbb{R}_{\ge0}\times\mathbb{R}^{n}\times\mathbb{R}^{n}$
and $q_{d}:\mathbb{R}_{\ge t_{0}}\to\mathbb{R}^{n}$, $\dot{q}_{d}:\mathbb{R}_{\ge t_{0}}\to\mathbb{R}^{n}$
and $\ddot{q}_{d}:\mathbb{R}_{\ge t_{0}}\to\mathbb{R}^{n}$ denote
generalized desired position, velocity, and acceleration, respectively.

The control objective is to design a controller such that $q$ converges
exponentially to $q_{d}$ despite the presence of non-vanishing disturbances,
unknown dynamics, and lack of knowledge of $\ddot{q}_{d}$. The subsequent
assumptions ensure feasibility of this objective.
\begin{assumption}
\label{asm:qd-bounds}\cite[Assumption 1]{Nino.Patil.ea2025} There
exist known constants $\overline{q}_{d},\overline{\dot{q}}_{d},\overline{\ddot{q}}_{d},\overline{\dddot{q}}_{d}\in\mathbb{R}_{\ge0}$
such that $\left\Vert q_{d}\left(t\right)\right\Vert \le\overline{q}_{d}$,
$\left\Vert \dot{q}_{d}\left(t\right)\right\Vert \le\overline{\dot{q}}_{d}$,
$\left\Vert \ddot{q}_{d}\left(t\right)\right\Vert \le\overline{\ddot{q}}_{d}$,
and $\left\Vert \dddot{q}_{d}\left(t\right)\right\Vert \le\overline{\dddot{q}}_{d}$
for all $t\in\mathbb{R}_{\ge0}$.
\end{assumption}
%
\begin{assumption}
\label{asm:d-bounds}\cite{Nino.Patil.ea2025} There exist known constants
$\overline{d},\overline{\dot{d}},\overline{\ddot{d}}\in\mathbb{R}_{\ge0}$
such that $\left\Vert d\left(t\right)\right\Vert \le\overline{d}$,
$\left\Vert \dot{d}\left(t\right)\right\Vert \le\overline{\dot{d}}$,
and $\left\Vert \ddot{d}\left(t\right)\right\Vert \le\overline{\ddot{d}}$
for all $t\in\mathbb{R}_{\ge0}$.\footnote{While adaptive RISE techniques like \cite{Bidikli.Tatlicioglu.ea2014a}
could be used to compensate for unknown bounds $\overline{d},\overline{\dot{d}},\overline{\ddot{d}}$,
the described formulation therein yields asymptotic convergence. Achieving
exponential convergence (as desired in this work) without Assumption
\ref{asm:d-bounds} remains an open challenge.}
\end{assumption}
%

\section{Control Design}

For subsequent analysis, let the tracking component of the controller
obtained from feedback linearization $u_{d}:\mathbb{R}_{\ge0}\to\mathbb{R}^{n}$
as
\begin{equation}
u_{d}\left(t\right)\triangleq d\left(t\right)+f\left(q_{d}\left(t\right),\dot{q}_{d}\left(t\right)\right)-\ddot{q}_{d}\left(t\right).\label{eq:u_d}
\end{equation}
Since $f$ is $\mathcal{C}^{1}$, Assumption \ref{asm:qd-bounds}
and extreme value theorem (EVT) guarantee that there exist constants
$\overline{f},\overline{\dot{f}}\in\mathbb{R}_{\ge0}$ such that $\left\Vert f\left(q_{d}\left(t\right),\dot{q}_{d}\left(t\right)\right)\right\Vert \le\overline{f}$
and $\left\Vert \dot{f}\left(q_{d}\left(t\right),\dot{q}_{d}\left(t\right)\right)\right\Vert \le\overline{\dot{f}}$
for all $t\in\mathbb{R}_{\ge0}$.
\begin{assumption}
\label{asm:ud_bounded} The constants $\overline{f},\overline{\dot{f}}\in\mathbb{R}_{\ge0}$
are known.
\end{assumption}
%
By Assumption \ref{asm:ud_bounded}, there exist known bounds $\overline{u}_{d}\triangleq\overline{d}+\overline{f}+\overline{\ddot{q}}_{d}$
and $\overline{\dot{u}}_{d}\triangleq\overline{\dddot{q}}_{d}+\overline{\dot{f}}\overline{\dot{q}}_{d}+\overline{\dot{f}}\overline{\ddot{q}}_{d}+\overline{\dot{d}}$
such that $\left\Vert u_{d}\left(t\right)\right\Vert \le\overline{u}_{d}$
and $\left\Vert \dot{u}_{d}\left(t\right)\right\Vert \le\overline{\dot{u}}_{d}$
for all $t\in\mathbb{R}_{\ge0}$. Subsequent analysis requires a
bound on the difference $\widetilde{u}\triangleq u_{d}-u$ in terms
of the ensemble system state $z:\mathcal{I}\to\mathbb{R}^{3n}$ defined
as $z\left(t\right)\triangleq\left[\begin{array}{ccc}
e^{\top}\left(t\right) & r_{1}^{\top}\left(t\right) & r_{2}^{\top}\left(t\right)\end{array}\right]^{\top}$. Since $f$ is $\mathcal{C}^{2}$, EVT, Assumption \ref{asm:qd-bounds},
and \cite[Lemma 4]{arxivKamalapurkar.Klotz.ea2013} guarantee that
there exist functions $\rho_{f},\rho_{f^{\prime}}:\mathbb{R}_{\ge0}\times\mathbb{R}_{\ge0}\to\mathbb{R}_{\ge0}$
that are class $\mathcal{P}_{\infty}$ with respect to each argument
such that $\left\Vert f\left(q,\dot{q}\right)-f\left(q_{d},\dot{q}_{d}\right)\right\Vert \le\rho_{f}\left(k_{1},\left\Vert z\right\Vert \right)\left\Vert z\right\Vert $
and $\left\Vert \left.\frac{\partial f}{\partial\dot{q}}\right|_{\left(q_{d},\dot{q}_{d}\right)}-\left.\frac{\partial f}{\partial\dot{q}}\right|_{\left(q,\dot{q}\right)}\right\Vert \le\rho_{f^{\prime}}\left(k_{1},\left\Vert z\right\Vert \right)\left\Vert z\right\Vert $.
\begin{assumption}
\label{asm:rho_f} The functions $\rho_{f}$ and $\rho_{f^{\prime}}$
are known.
\end{assumption}
%
Subtracting (\ref{eq:u_d}) from (\ref{eq:agent-dynamics}) and using
(\ref{eq:r1_def}), (\ref{eq:r2_def}), and the time derivative of
(\ref{eq:r2_def}) gives that
\begin{equation}
\begin{alignedat}{1}\widetilde{u}\left(t\right) & =r_{2}\left(t\right)-k_{1}\dot{e}\left(t\right)-k_{2}r_{1}\left(t\right)\\
 & -f\left(q_{d}\left(t\right),\dot{q}_{d}\left(t\right)\right)+f\left(q\left(t\right),\dot{q}\left(t\right)\right),
\end{alignedat}
\label{eq:u_tilde}
\end{equation}
for all $t\in\mathcal{I}$. By Assumption \ref{asm:rho_f}, there
exists a known function $\rho_{\widetilde{u}}:\mathbb{R}_{\ge0}\times\mathbb{R}_{\ge0}\times\mathbb{R}_{\ge0}\to\mathbb{R}_{\ge0}$that
is of class $\mathcal{P}_{\infty}$ with respect to each argument
such that $\left\Vert \widetilde{u}\right\Vert \le\rho_{\widetilde{u}}\left(k_{1},k_{2},\left\Vert z\right\Vert \right)\left\Vert z\right\Vert $
where $\rho_{\widetilde{u}}\left(k_{1},k_{2},s\right)\triangleq1+k_{1}+k_{1}^{2}+k_{2}+\rho_{f}\left(k_{1},s\right)$.

To facilitate subsequent control design, the first and second auxiliary
tracking errors are defined as
\begin{align}
r_{1}\left(t\right) & \triangleq\dot{e}\left(t\right)+k_{1}e\left(t\right),\label{eq:r1_def}\\
r_{2}\left(t\right) & \triangleq\dot{r}_{1}\left(t\right)+k_{2}r_{1}\left(t\right),\label{eq:r2_def}
\end{align}
where $k_{1},k_{2}\in\mathbb{R}_{>0}$ are user-selected gains. Differentiating
(\ref{eq:r1_def}) and using (\ref{eq:agent-dynamics}), (\ref{eq:reference-dynamics}),
and ($\ref{eq:r1_def}$) gives
\begin{equation}
\begin{alignedat}{1}\dot{r}_{1}\left(t\right) & =h\left(q_{d}\left(t\right),\dot{q}_{d}\left(t\right)\right)-f\left(q\left(t\right),\dot{q}\left(t\right)\right)\\
 & -u\left(t\right)-d\left(t\right)+k_{1}r_{1}\left(t\right)-k_{1}^{2}e\left(t\right),
\end{alignedat}
\label{eq:r1-step-1}
\end{equation}
for all $t\in\mathcal{I}$. Differentiating (\ref{eq:r2_def}) and
substituting (\ref{eq:r1-step-1}) and ($\ref{eq:r1_def}$) gives
\begin{equation}
\begin{alignedat}{1}\dot{r}_{2}\left(t\right) & =F\left(\kappa\left(t\right)\right)-\left.\frac{\partial f}{\partial\dot{q}}\right|_{\left(q\left(t\right),\dot{q}\left(t\right)\right)}d\left(t\right)-\dot{u}\left(t\right)-\dot{d}\left(t\right)\\
 & +\left(k_{1}+k_{2}\right)r_{2}\left(t\right)-\left(k_{1}k_{2}+k_{2}^{2}+k_{1}^{2}\right)r_{1}\left(t\right)+k_{1}^{3}e\left(t\right),
\end{alignedat}
\label{eq:r2-step-1}
\end{equation}
for all $t\in\mathcal{I}$, where $\kappa:\mathcal{I}\to\mathbb{R}^{5n}$
is defined as $\kappa\triangleq\left[\begin{array}{ccccc}
q^{\top} & \dot{q}^{\top} & q_{d}^{\top} & \dot{q}_{d}^{\top} & u^{\top}\end{array}\right]^{\top}$ and $F:\mathbb{R}^{5n}\to\mathbb{R}^{n}$ is defined as $F\left(\kappa\right)\triangleq\left.\frac{\partial h}{\partial q_{d}}\right|_{\left(q_{d},\dot{q}_{d}\right)}\dot{q}_{d}+\left.\frac{\partial h}{\partial\dot{q}_{d}}\right|_{\left(q_{d},\dot{q}_{d}\right)}h\left(q_{d},\dot{q}_{d}\right)-\left.\frac{\partial f}{\partial q}\right|_{\left(q,\dot{q}\right)}\dot{q}-\left.\frac{\partial f}{\partial\dot{q}}\right|_{\left(q,\dot{q}\right)}f\left(q,\dot{q}\right)-\left.\frac{\partial f}{\partial\dot{q}}\right|_{\left(q,\dot{q}\right)}u$.

\subsection{Function Approximation}

Since $f$ and $h$ are unknown, $F$ is unknown and will be approximated
online with a user-selected parametric universal approximator as $\Phi\left(\kappa,\widehat{\theta}\right)$
where $\Phi:\mathbb{R}^{L_{\text{in}}}\times\mathbb{R}^{{\tt p}}\to\mathbb{R}^{L_{\text{out}}}$
is a function approximator architecture mapping of class $\mathcal{C}^{2}$,
$\widehat{\theta}\in\mathbb{R}^{{\tt p}}$ is a parameter, $L_{\text{out}}=n$,
$L_{\text{in}}=5n$, and ${\tt p}\in\mathbb{Z}_{\ge1}$ is the number
of parameters.

\begin{rem}
In this work, any class $\mathcal{C}^{2}$ parametric approximator
which satisfies a universal approximation theorem like Theorem \ref{thm:Universal-Approximation-Theorem}
can be used for $\Phi$, in which there must exist a sufficiently
sized architecture (e.g., network width for shallow neural networks,
network depth for deep neural networks of fixed width, or number of
blocks for ResNets) and a corresponding ideal parameter $\theta^{\ast}\in\mathbb{R}^{{\tt p}}$
such that any function of class $\mathcal{C}^{1}$ can be approximated
to arbitrarily small precision under the $C^{1}$ norm on a compact
set $\Omega\subsetneq\mathbb{R}^{L_{\text{in}}}$. For example, radial
basis functions, Chebyshev polynomials, transformers, or Kolmogorov-Arnold
networks can be used.
\end{rem}
%
To construct the compact domain of inputs over which the approximation
of $F$ by $\Phi\left(\kappa,\widehat{\theta}\right)$ is of interest,
the idealized input $\kappa_{d}:\mathbb{R}_{\ge0}\to\mathbb{R}^{5n}$
seen when (\ref{eq:agent-dynamics}) perfectly tracks (\ref{eq:reference-dynamics})
is defined as $\kappa_{d}\triangleq\left[\begin{array}{ccccc}
q_{d}^{\top} & \dot{q}_{d}^{\top} & q_{d}^{\top} & \dot{q}_{d}^{\top} & u_{d}^{\top}\end{array}\right]^{\top}\in\mathbb{R}^{5n}$. By Assumptions \ref{asm:qd-bounds} and \ref{asm:ud_bounded}, there
exist known constants $\overline{\kappa}_{d}\triangleq2\overline{q}_{d}+2\overline{\dot{q}}_{d}+\overline{u}_{d}$
and $\overline{\dot{\kappa}}_{d}\triangleq2\overline{\dot{q}}_{d}+2\overline{\ddot{q}}_{d}+\overline{\dot{u}}_{d}$
such that $\left\Vert \kappa_{d}\right\Vert \le\overline{\kappa}_{d}$
and $\left\Vert \dot{\kappa}_{d}\right\Vert \le\overline{\dot{\kappa}}_{d}$
for all $t\in\mathbb{R}_{\ge0}$. Using ($\ref{eq:r1_def}$), Assumption
\ref{asm:rho_f}, and the definition of $\rho_{\widetilde{u}}$, there
exists a known function that is of class $\mathcal{P}_{\infty}$ with
respect to each argument such that $\left\Vert \widetilde{\kappa}\right\Vert \le\rho_{\widetilde{\kappa}}\left(k_{1},k_{2},\left\Vert z\right\Vert \right)\left\Vert z\right\Vert $,
where $\widetilde{\kappa}\triangleq\kappa_{d}-\kappa$ and $\rho_{\widetilde{\kappa}}\left(k_{1},k_{2},s\right)\triangleq k_{1}+2+\rho_{\widetilde{u}}\left(k_{1},k_{2},s\right)$.
The approximation domain of interest for all inputs $\kappa$ is given
as
\begin{equation}
\Omega\triangleq\left\{ \zeta\in\mathbb{R}^{5n}:\left\Vert \zeta\right\Vert \le\overline{\kappa}\right\} ,\quad\overline{\kappa}\triangleq\overline{\kappa}_{d}+\rho_{\widetilde{\kappa}}\left(k_{1},k_{2},C_{\Delta}\right)C_{\Delta},\label{eq:Omega}
\end{equation}
where $C_{\Delta}\in\mathbb{R}_{>0}$ is a subsequently defined constant.

The parameter $\widehat{\theta}$ is intended to dynamically estimate
$\theta^{\ast}$ and evolve continuously but be constrained to a user-defined
bounded search space $\mho$. Without loss of generality, it is assumed
herein that $\mho=\overline{B}_{\overline{\theta}}\left(0_{{\tt p}}\right)$,
where $\overline{\theta}\in\mathbb{R}_{>0}$ is a user-defined radius.
The rule governing the trajectory of $\widehat{\theta}$ over time
is called the update law and is chosen by the user. The update law
must satisfy the following properties.
\begin{property}
\label{prop:theta_bar}\cite[Assumption 1]{Patil.Isaly.ea2022} The
parameter's trajectory $\widehat{\theta}:\mathbb{R}_{\ge0}\to\mathbb{R}^{{\tt p}}$
is of class $\mathcal{C}^{1}$ and there exists known constants $\overline{\theta},\eta_{1}\in\mathbb{R}_{>0}$
such that the trajectory and its first time derivative $\dot{\widehat{\theta}}:\mathbb{R}_{\ge0}\to\mathbb{R}^{{\tt p}}$
satisfy $\left\Vert \widehat{\theta}\left(t\right)\right\Vert \le\overline{\theta}$
and $\left\Vert \dot{\widehat{\theta}}\left(t\right)\right\Vert \le\eta_{1}$
for all $t\in\mathbb{R}_{\ge0}$.
\end{property}
%
\begin{property}
\label{prop:weight-initialization} There exists a known constant
$\overline{\varepsilon}_{\text{d}}\in\mathbb{R}_{>0}$ such that parameter's
trajectory satisfies $E_{\Omega}\left(\widehat{\theta}\left(t\right)\right)\le\overline{\varepsilon}_{\text{d}}$
for all $t\in\mathbb{R}_{\ge0}$, where $\varepsilon\left(\kappa,\theta\right)\triangleq\Phi\left(\kappa,\theta\right)-F\left(\kappa\right)$
and $E_{\Omega}\left(\theta\right)\triangleq\sup_{\kappa\in\Omega}\left\Vert \varepsilon\left(\kappa,\theta\right)\right\Vert +\sup_{\kappa\in\Omega}\left\Vert \left.\frac{d\varepsilon}{d\kappa}\right|_{\left(\kappa,\theta\right)}\right\Vert $.
\end{property}
%
Property \ref{prop:theta_bar} can be met, for instance, by using
a projection algorithm and limiting the speed of parameter changes.
Satisfying Property \ref{prop:weight-initialization} requires careful
consideration of both parameter initialization and the update law.
The initial parameter $\widehat{\theta}\left(t_{0}\right)$ must either
be obtained by pre-training the function approximator on input-output
pairs generated by a function representative of $F$ or by using the
following approach. Notice that EVT ensures that $\left\Vert F\right\Vert +\left\Vert \frac{dF}{d\kappa}\right\Vert $
achieves a finite maximum over $\Omega$ since $\Omega$ is compact
and $F$ is of class $\mathcal{C}^{1}$. Assume a bound $\overline{F}\in\mathbb{R}_{\ge0}$
on such a maximum is known. Pick $\widehat{\theta}\left(t_{0}\right)$
and a suitable $\mho$ such that $\Phi\left(\kappa,\widehat{\theta}\left(t_{0}\right)\right)=0_{n}$
for all $\kappa\in\Omega$. Finally, select $\overline{\varepsilon}_{\text{d}}=\overline{F}$.
Satisfying Property \ref{prop:weight-initialization} beyond the initial
time requires that any information gained online to refine the parameter
not result in a $E_{\Omega}\left(\widehat{\theta}\left(t\right)\right)$
that exceeds $\overline{\varepsilon}_{\text{d}}$. Heuristically,
this phenomena, called catastrophic forgetting, can be avoided by
using experience replay methods like concurrent learning (CL).\footnote{For example, a first-order CL-based update law for parameters in universal
approximators like $\Phi$ can be found in \cite{Hart.Patil.ea2025}.
Note that any discrete updates to a history stack made online may
not satisfy the $\mathcal{C}^{1}$ requirement of Property \ref{prop:theta_bar}.
However, online history stack management is not required to prevent
catastrophic forgetting.}

Since the expressive powers of ResNets and many similar approximators
rely on the ability to select $\theta^{\ast}$ from the whole of $\mathbb{R}^{{\tt p}}$
(cf. Theorem \ref{thm:Universal-Approximation-Theorem}) rather than
a prescribed $\mho$, $F$ must belong to a strict subset of all continuous
functions as a consequence of Property \ref{prop:theta_bar}.
\begin{assumption}
\label{asm:bounded-weights} The function $F$ and its first derivative
$\frac{dF}{d\kappa}$ can be uniformly approximated to arbitrary precision
using bounded weights for a sufficiently large approximator architecture.
Formally, $F\in\{\mathcal{C}^{1}(\Omega;\mathbb{R}^{n}):\forall\overline{\varepsilon}\in\mathbb{R}_{>0},\exists{\tt p}\in\mathbb{Z}_{\ge1}:\exists\mho\subsetneq\mathbb{R}^{{\tt p}}:\exists\theta^{\ast}\in\mho:E_{\Omega}(\theta^{\ast})\le\overline{\varepsilon}\}$.\footnote{Assumption \ref{asm:bounded-weights} is simply a formal statement
of an implicit assumption common to many previous adaptive controls
works using parametric approximators.}
\end{assumption}
%
Assumption \ref{asm:bounded-weights} is required because of the following
tradeoff. In parametric nonlinear regression, the user is forced to
pick two of the following three desirable qualities for their approximator:
(i) $F$ can be any $\mathcal{C}^{1}$ function, (ii) $F$ can be
approximated arbitrarily well ($\overline{\varepsilon}_{\text{d}}\to0$
as ${\tt p}\to\infty$), (iii) a known bound exists on the weights
($\overline{\theta}$) needed to achieve the desired accuracy for
the target function. Without (ii), there will exist sufficiently small
values of $\overline{\varepsilon}_{\text{d}}$ that are not achievable
for the given $\mho$ even as ${\tt p}\to\infty$. That is, the infimum
of all values $\overline{\varepsilon}_{\text{d}}$ may achieve is
an unknown positive value (rather than zero) that might be arbitrarily
large. Thus, (ii) is required. Since computers possess finite memory,
(iii) is required. Thus, (i) must be sacrificed in place of Assumption
\ref{asm:bounded-weights}.

By Assumption \ref{asm:bounded-weights}, the set $\Theta^{\ast}\triangleq\arg\min_{\vartheta\in\mho}E_{\Omega}\left(\theta\right)$
is nonempty and all $\theta^{\ast}\in\Theta^{\ast}$ satisfy $\left\Vert \theta^{\ast}\right\Vert \le\overline{\theta}$.
Since the design of the update law and the choice of parametric approximator
for $\Phi$ are at the user's discretion, the approach in this work
is said to be modular.

\subsection{Control Law}

Motivated by the form of (\ref{eq:r2-step-1}), the controller is
designed as
\begin{equation}
\begin{alignedat}{1}\dot{u}\left(t\right) & =\left(k_{1}+k_{2}+k_{3}\right)r_{2}\left(t\right)+\left(1-k_{1}^{2}-k_{2}^{2}-k_{1}k_{2}\right)r_{1}\left(t\right)\\
 & +k_{1}^{3}e\left(t\right)+K_{\text{RISE}}\text{sgn}\left(r_{1}\left(t\right)\right)+\Phi\left(\kappa\left(t\right),\hat{\theta}\left(t\right)\right),
\end{alignedat}
\label{eq:u_dot}
\end{equation}
where $K_{\text{RISE}},k_{3}\in\mathbb{R}_{>0}$ are user-defined
constants. Integrating both sides of (\ref{eq:u_dot}) and picking
the initial condition $u\left(t_{0}\right)=K_{P}e\left(t_{0}\right)+K_{D}\dot{e}\left(t_{0}\right)$
yields the implementable form
\begin{equation}
\begin{alignedat}{1}u\left(t\right) & =K_{P}e\left(t\right)+K_{D}\dot{e}\left(t\right)+\intop_{t_{0}}^{t}\left(K_{I}e\left(\tau\right)\right)d\tau\\
 & +\left(\intop_{t_{0}}^{t}K_{\text{RISE}}\text{sgn}\left(r_{1}\left(\tau\right)\right)+\Phi\left(\kappa\left(\tau\right),\hat{\theta}\left(\tau\right)\right)\right)d\tau
\end{alignedat}
\label{eq:u_implementable}
\end{equation}
\[
\]
where $K_{P}\triangleq k_{1}k_{2}+k_{1}k_{3}+k_{2}k_{3}+1$, $K_{I}\triangleq k_{1}k_{2}k_{3}+k_{1}$,
and $K_{D}\triangleq k_{1}+k_{2}+k_{3}$.\footnote{Although $\ddot{q}$ is not assumed to be measurable, it appears in
(\ref{eq:u_dot}) by virtue of $r_{2}$. However, note that the implementable
form only requires $q$ and $\dot{q}$. } Substituting (\ref{eq:u_dot}) and the definition of $\varepsilon$
into (\ref{eq:r2-step-1}) and performing algebraic manipulation gives
that $r_{2}:\mathcal{I}\to\mathbb{R}^{n}$ is a Filippov solution
to
\[
\dot{r}_{2}\left(t\right)\in\widetilde{N}\left(\kappa,\kappa_{d},\theta^{\ast},\hat{\theta},t\right)-k_{3}r_{2}\left(t\right)-r_{1}\left(t\right)-K_{\text{RISE}}K[\text{sgn}]\left(r_{1}\left(t\right)\right)+D\left(\kappa_{d},\theta^{\ast},\hat{\theta},t\right),
\]
where
\begin{align*}
D\left(\kappa_{d},\theta^{\ast},\hat{\theta},t\right) & \triangleq-\dot{d}\left(t\right)-\left.\frac{\partial f}{\partial\dot{q}}\right|_{\left(q_{d},\dot{q}_{d}\right)}d\left(t\right)+\varepsilon\left(\kappa_{d}\right)+N_{0}\left(\kappa_{d},\theta^{\ast},\hat{\theta}\right),\\
N_{0}\left(\kappa,\theta^{\ast},\hat{\theta}\right) & \triangleq\Phi\left(\kappa,\theta^{\ast}\right)-\Phi\left(\kappa,\hat{\theta}\right),\\
\widetilde{N}\left(\kappa,\kappa_{d},\theta^{\ast},\hat{\theta},t\right) & \triangleq\widetilde{N}_{1}\left(\kappa,\kappa_{d},\theta^{\ast},\hat{\theta}\right)+\widetilde{N}_{\varepsilon}\left(\kappa,\kappa_{d}\right)+\widetilde{N}_{2}\left(q,\dot{q},q_{d},\dot{q}_{d}\right)d\left(t\right),\\
\widetilde{N}_{\varepsilon}\left(\kappa,\kappa_{d}\right) & \triangleq\varepsilon\left(\kappa\right)-\varepsilon\left(\kappa_{d}\right),\\
\widetilde{N}_{1}\left(\kappa,\kappa_{d},\theta^{\ast},\hat{\theta}\right) & \triangleq N_{0}\left(\kappa,\theta^{\ast},\hat{\theta}\right)-N_{0}\left(\kappa_{d},\theta^{\ast},\hat{\theta}\right),\\
\widetilde{N}_{2}\left(q,\dot{q},q_{d},\dot{q}_{d}\right) & \triangleq\left.\frac{\partial f}{\partial\dot{q}}\right|_{\left(q_{d},\dot{q}_{d}\right)}-\left.\frac{\partial f}{\partial\dot{q}}\right|_{\left(q,\dot{q}\right)}
\end{align*}
Since $\Phi$ is of class $\mathcal{C}^{1}$ and $\kappa_{d}$ is
bounded, using \cite[Lemma 4]{arxivKamalapurkar.Klotz.ea2013}, the
bound on $\left\Vert \widetilde{\kappa}\right\Vert $ and Property
\ref{asm:bounded-weights}, there exists a known function $\rho_{\Phi}:\mathbb{R}_{\ge0}\times\mathbb{R}_{\ge0}\times\mathbb{R}_{\ge0}\to\mathbb{R}_{\ge0}$
that is of class $\mathcal{P}_{\infty}$ with respect to each argument
such that $\left\Vert \widetilde{N}_{1}\left(\kappa,\kappa_{d},\theta^{\ast},\hat{\theta}\right)\right\Vert \le\rho_{\Phi}\left(k_{1},k_{2},\left\Vert z\right\Vert \right)\left\Vert z\right\Vert $.
Since $\overline{\kappa}_{d}\le\overline{\kappa}$, Property \ref{prop:weight-initialization}
ensures that $\left\Vert \varepsilon\left(\kappa_{d}\left(t\right)\right)\right\Vert \le\overline{\varepsilon}_{\text{d}}$
and $\left\Vert \left.\frac{d\varepsilon}{d\kappa}\right|_{\kappa_{d}\left(t\right)}\right\Vert \le\overline{\varepsilon}_{\text{d}}$
for all $t\in\mathbb{R}_{\ge0}$. Since $\left\Vert \kappa\right\Vert \le\overline{\kappa}_{d}+\left\Vert \widetilde{\kappa}\right\Vert $,
the bound on $\left\Vert \widetilde{\kappa}\right\Vert $ gives that
$\left\Vert \kappa\right\Vert \le\overline{\kappa}$ for all $\left\Vert z\right\Vert \le C_{\Delta}$.
Since $\varepsilon$ is of class $\mathcal{C}^{1}$, it follows that
$\left\Vert \widetilde{N}_{\varepsilon}\left(\kappa,\kappa_{d}\right)\right\Vert \le\overline{\varepsilon}_{\text{d}}\rho_{\widetilde{\kappa}}\left(k_{1},k_{2},\left\Vert z\right\Vert \right)\left\Vert z\right\Vert $
for all $\left\Vert z\right\Vert \le C_{\Delta}$ \cite[Lemma 4]{arxivKamalapurkar.Klotz.ea2013}.
Since $f$ is of class $\mathcal{C}^{2}$, Assumption \ref{asm:qd-bounds}
and EVT guarantee that there exist constants $\overline{f^{\prime}},\overline{f^{\prime\prime}}\in\mathbb{R}_{\ge0}$
such $\left\Vert \left.\frac{\partial f}{\partial\dot{q}}\right|_{\left(q_{d}\left(t\right),\dot{q}_{d}\left(t\right)\right)}\right\Vert \le\overline{f^{\prime}}$,
$\left\Vert \left.\frac{\partial f}{\partial\dot{q}\partial q}\right|_{\left(q_{d},\dot{q}_{d}\right)}\right\Vert \le\overline{f^{\prime\prime}}$,
and $\left\Vert \left.\frac{\partial^{2}f}{\partial\dot{q}^{2}}\right|_{\left(q_{d},\dot{q}_{d}\right)}\right\Vert \le\overline{f^{\prime\prime}}$
for all $t\in\mathbb{R}_{\ge0}$. 
\begin{assumption}
\label{asm:f-prime} The constants $\overline{f^{\prime}}$ and $\overline{f^{\prime\prime}}$
are known.
\end{assumption}
%
Since $\Phi$ is of class $\mathcal{C}^{0}$ and $\left\Vert \kappa_{d}\left(t\right)\right\Vert $,
$\left\Vert \dot{\kappa}_{d}\left(t\right)\right\Vert $, $\left\Vert \varepsilon\left(\kappa_{d}\left(t\right)\right)\right\Vert $,
and $\left\Vert \left.\frac{d\varepsilon}{d\kappa}\right|_{\kappa_{d}\left(t\right)}\right\Vert $
are bounded by known constants, EVT, Property \ref{asm:bounded-weights},
and Assumptions \ref{asm:d-bounds} and \ref{asm:f-prime} ensure
that there exist known constants $\overline{D},\overline{\dot{D}}\in\mathbb{R}_{\ge0}$
such that $\left\Vert D\left(\kappa_{d},\theta^{\ast},\hat{\theta},t\right)\right\Vert \le\overline{D}$
and $\left\Vert \dot{D}\left(\kappa_{d},\theta^{\ast},\hat{\theta},t\right)\right\Vert \le\overline{\dot{D}}$
for all $t\in\mathbb{R}_{\ge0}$. By Assumption \ref{asm:rho_f},
$\left\Vert \widetilde{N}_{2}\left(q,\dot{q},q_{d},\dot{q}_{d}\right)\right\Vert \le\rho_{f^{\prime}}\left(k_{1},\left\Vert z\right\Vert \right)\left\Vert z\right\Vert $.
Using Assumption \ref{asm:d-bounds}, it follows that there exists
a known function $\rho_{0}:\mathbb{R}_{\ge0}\times\mathbb{R}_{\ge0}\times\mathbb{R}_{\ge0}\to\mathbb{R}_{\ge0}$
that is of class $\mathcal{P}_{\infty}$ with respect to each argument
such that $\left\Vert \widetilde{N}\left(\kappa,\kappa_{d},\theta^{\ast},\hat{\theta},t\right)\right\Vert \le\rho_{0}\left(k_{1},k_{2},\left\Vert z\right\Vert \right)\left\Vert z\right\Vert $
for all $\left\Vert z\right\Vert \le C_{\Delta}$, where $\rho_{0}\left(k_{1},k_{2},s\right)\triangleq\rho_{\Phi}\left(k_{1},k_{2},s\right)+\overline{\varepsilon}_{\text{d}}\rho_{\widetilde{\kappa}}\left(k_{1},k_{2},s\right)+\overline{d}\rho_{f^{\prime}}\left(k_{1},s\right)$.

\section{Stability Analysis}

\bibliographystyle{ieeetr}
\bibliography{bib/sources,\string"/Users/max/Library/CloudStorage/SynologyDrive-1/1 - Academics/1 - University of Florida/1 - Research/ncrbibs/master\string",\string"/Users/max/Library/CloudStorage/SynologyDrive-1/1 - Academics/1 - University of Florida/1 - Research/ncrbibs/ncr\string",\string"/Users/max/Library/CloudStorage/SynologyDrive-1/1 - Academics/1 - University of Florida/1 - Research/ncrbibs/encr\string"}

\end{document}
